\newenvironment{warning}
 {\par\begin{mdframed}[linewidth=2pt,linecolor=red]%
    \begin{list}{}{\leftmargin=1cm
                   \labelwidth=\leftmargin}}
  {\end{list}\end{mdframed}\par}

\subchapter{Setup AHAB on i.MX93}{Objective: Learn how to implement the first stage of a secure boot}

\section{Prerequisite}

To build the AHAB container for the first stage of the i.MX93 secure boot, we are going
to use NXP's Secure Provisioning SDK (\href{https://spsdk.readthedocs.io/en/latest/}{SPSDK}).

First, install it in your python virtual environment (the one you created to install kas
in the \code{Build Yocto} section of the \code{Board setup} lab):

\begin{bashinput}
(.venv) $ pip install spsdk
\end{bashinput}

AHAB also necessitates asymmetric keys to:
\begin{itemize}
  \item sign the container on the host
  \item verify the signature on the target
\end{itemize}

It can include up to 4 of these keys, which must all be hashed together, and we therefore
need to generate those 4 Super Root Keys (SRKs):

\begin{bashinput}
$ for i in `seq 0 3`; \
  do openssl genrsa -out keys/srk_$i.pem 4096; \
  openssl rsa -in keys/srk_$i.pem -pubout > keys/srk_$i.pub; \
done
\end{bashinput}

\section{Building a signed AHAB container}

The first step in building the AHAB image container we will use for the first stage of secure
boot, is to identify the system's components we want the ROM code to verify.

Remember that the container will actually be split into 2 container sets:
\begin{itemize}
  \item The \textbf{primary} container set, which will contain the very early boot components
  \item The \textbf{secondary} container set, containing the next stage(s)
\end{itemize}

Both will be signed with one of the SRKs created at the previous step.

The \textbf{primary} container set will contain:
\begin{itemize}
  \item The ELE firmware, which must be downloaded from NXP. In our case, Yocto did that for us:
	  \code{build/tmp-glibc/sysroots-components/cortexa55/firmware-ele-imx/usr/lib/firmware/imx/ele/mx93a1-ahab-container.img}
  \item The DRAM training binaries, which Yocto also downloaded:
  \begin{itemize}
	  \item \code{build/tmp-glibc/sysroots-components/cortexa55/imx-boot-firmware-files/firmware/lpddr4_imem_1d_v202201.bin}
	  \item \code{build/tmp-glibc/sysroots-components/cortexa55/imx-boot-firmware-files/firmware/lpddr4_imem_2d_v202201.bin}
	  \item \code{build/tmp-glibc/sysroots-components/cortexa55/imx-boot-firmware-files/firmware/lpddr4_dmem_1d_v202201.bin}
	  \item \code{build/tmp-glibc/sysroots-components/cortexa55/imx-boot-firmware-files/firmware/lpddr4_dmem_2d_v202201.bin}
  \end{itemize}
  \item Finally, U-Boot's SPL built by Yocto, and located in \code{build/tmp-glibc/work/freiheit93-oe-linux/u-boot-kiss/git/build/imx93_frdm_defconfig/u-boot-spl-ddr.bin}
\end{itemize}

The \textbf{secondary} container set will contain:
\begin{itemize}
  \item TF-A's BL31: \code{build/tmp-glibc/deploy/images/freiheit93/bl31.bin}
  \item OP-TEE: \code{build/tmp-glibc/deploy/images/freiheit93/optee/tee-raw.bin}
  \item Finally, U-Boot proper: \code{build/tmp-glibc/work/freiheit93-oe-linux/u-boot-kiss/git/build/imx93_frdm_defconfig/u-boot-dtb.bin}
\end{itemize}

The easiest way to write the YAML templates is to start from the ones included in spsdk:

\begin{bashinput}
(.venv) $ nxpimage bootable-image get-templates -f mimx9352 -t imx_boot_flash_singleboot -o ahab_template
(.venv) $ ls ahab_template
bootable_image.yaml  inputs  spl.yaml  uboot.yaml
\end{bashinput}

Once you are done filling the YAML templates, you can build the image:

\begin{bashinput}
(.venv) $ nxpimage -v bootable-image export -c ahab_template/bootable_image.yaml
\end{bashinput}

This should create:
\begin{itemize}
  \item The AHAB container image: \code{ahab_template/bootable_image.bin} 
  \item One \code{nxpele} script for each subcontainer:
  \begin{itemize}
    \item \code{output/ahab_oem0_srk0_hash_nxpele.bcf}
    \item \code{output/ahab_oem1_srk0_hash_nxpele.bcf}
  \end{itemize}	
\end{itemize}

The SoC only has one set of SRK eFuses, so both container set should use the same keys,
so these two scripts should be identical:

\begin{bashinput}
$ diff output/ahab_oem1_srk0_hash_nxpele.bcf output/ahab_oem0_srk0_hash_nxpele.bcf
\end{bashinput}

\section{Flashing the i.MX93 eFuses}

\begin{warning}
\Large
This step is irreversible: the eFuses can only be flashed once. Although the device will still
boot unless the i.MX93 is transitioned into \textbf{OEM\_CLOSED} state, this might still be undesirable.
\end{warning}

The easiest way to interact with the SoC's eFuses is by using \href{https://spsdk.readthedocs.io/en/latest/}{SPSDK}'s
\code{nxpele} tool. In our case, we will use it to issue command to U-Boot over the command line.
We therefore need to reset the board, and interrupt U-Boot's boot process.

Once this is done, we'll need to disconnect the serial client, as it would interfere with \code{nxpele}.

\begin{bashinput}
$ for i in `seq 128 135`; do nxpele -p /dev/ttyACM0 -d uboot_serial -f mimx9352 -v read-common-fuse --index $i; done
INFO:spsdk.ele.ele_comm:ELE communicator is using 131072 B size buffer at 83800000 address in mimx9352, Revision: latest target.
INFO:spsdk.ele.ele_comm:Sent message information:
Command:         READ_COMMON_FUSE - (0x97)
Command words:   2
Command data:    False
Response words:  3
Response data:   False
Response status: Success

Read common fuse ends successfully.
Fuse ID_128: 0x00000000

[...]

Read common fuse ends successfully.
Fuse ID_135: 0x00000000
\end{bashinput}

Before you flash the eFuses, take note of the output of U-Boot's \code{ahab_status}:

If the board is not booting from a \textbf{signed} AHAB container, the output should be:
\begin{bashinput}
$ Hit any key to stop autoboot:  0 
u-boot=> ahab_status
Lifecycle: 0x00000008, OEM Open


        0x0287eed6
        IPC = MU APD (0x2)
        CMD = ELE_OEM_CNTN_AUTH_REQ (0x87)
        IND = ELE_NO_AUTHENTICATION_FAILURE_IND (0xEE)
        STA = ELE_SUCCESS_IND (0xD6)

        0x0287eed6
        IPC = MU APD (0x2)
        CMD = ELE_OEM_CNTN_AUTH_REQ (0x87)
        IND = ELE_NO_AUTHENTICATION_FAILURE_IND (0xEE)
        STA = ELE_SUCCESS_IND (0xD6)

\end{bashinput}

If the board \textbf{is} booting from a \textbf{signed} AHAB container, this should be the output:

\begin{bashinput}
Hit any key to stop autoboot:  0 
u-boot=> ahab_status
Lifecycle: 0x00000008, OEM Open


        0x0287fad6
        IPC = MU APD (0x2)
        CMD = ELE_OEM_CNTN_AUTH_REQ (0x87)
        IND = ELE_BAD_KEY_HASH_FAILURE_IND (0xFA)
        STA = ELE_SUCCESS_IND (0xD6)

        0x0287fad6
        IPC = MU APD (0x2)
        CMD = ELE_OEM_CNTN_AUTH_REQ (0x87)
        IND = ELE_BAD_KEY_HASH_FAILURE_IND (0xFA)
        STA = ELE_SUCCESS_IND (0xD6)

\end{bashinput}

We can now flash the eFuses. Although this step is \textbf{irreversible}, it will not prevent your board
from booting, as AHAB will not enforce the signature in \textbf{OEM\_OPEN}.

\begin{bashinput}
$ nxpele -p /dev/ttyACM0 -d uboot_serial -f mimx9352 -v batch output/ahab_oem0_srk0_hash_nxpele.bcf 
INFO:spsdk.ele.ele_comm:ELE communicator is using 131072 B size buffer at 83800000 address in mimx9352, Revision: latest target.
INFO:spsdk.ele.ele_comm:Sent message information:
Command:         WRITE_FUSE - (0xd6)
Command words:   3
Command data:    False
Response words:  3
Response data:   False
Response status: Success

ELE write fuse ends successfully.

[...] Repeats 7 more times because the hash is 8*32 = 256 bits

\end{bashinput}

We can now re-read the fuses, and they should contain our hash:

\begin{bashinput}
$ for i in `seq 128 135`; do nxpele -p /dev/ttyACM0 -d uboot_serial -f mimx9352 -v read-common-fuse --index $i; done
INFO:spsdk.ele.ele_comm:ELE communicator is using 131072 B size buffer at 83800000 address in mimx9352, Revision: latest target.
INFO:spsdk.ele.ele_comm:Sent message information:
Command:         READ_COMMON_FUSE - (0x97)
Command words:   2
Command data:    False
Response words:  3
Response data:   False
Response status: Success

Read common fuse ends successfully.
Fuse ID_128: 0x09C79D69

[...]
Fuse ID_129: 0x340DBD87
[...]
Fuse ID_130: 0x6BA8EB5A
[...]
Fuse ID_131: 0xECD130AD
[...]
Fuse ID_132: 0xFD09593E
[...]
Fuse ID_133: 0x38FA50F1
[...]
Fuse ID_134: 0x18DC2FA2
[...]
Fuse ID_135: 0xC6DED0B6
\end{bashinput}

\section{Testing the AHAB stage}

Now that the eFuses are flashed (if you finished the previous step), the ELE ROM should now be able
to verify the signature on the AHAB subcontainers.

If you succesfully signed your AHAB containers, this is the output you should see when running \code{ahab_status}:

\begin{bashinput}
Hit any key to stop autoboot:  0 
u-boot=> ahab_status
Lifecycle: 0x00000008, OEM Open


        No Events Found!

\end{bashinput}

This indicates that the system would have booted, even in \code{OEM_CLOSED} or \code{OEM_LOCKED}.

Because we have not transitioned the system to either state, though, it will still boot an unsigned
container, or a wrongly signed one.

For example, we can change the value of \code{used_srk_id} in \code{ahab_template/uboot.yaml} to \code{1}:
\begin{bashinput}
u-boot=> ahab_status
Lifecycle: 0x00000008, OEM Open


        0x0287f0d6
        IPC = MU APD (0x2)
        CMD = ELE_OEM_CNTN_AUTH_REQ (0x87)
        IND = ELE_BAD_SIGNATURE_FAILURE_IND (0xF0)
        STA = ELE_SUCCESS_IND (0xD6)

\end{bashinput}

Notice that we only have one error this time, because we only changed one subcontainer.
We also have a very slightly different error code.

To get back to good state, you can either change \code{used_srk_id} back to \code{0}, or adapt
\code{signer} to use the SRK with ID \code{1}.
